\documentclass[12pt, xcolor={usenames,dvipsnames,svgnames,x11names,table}]{beamer}

%%%%% make text revisible when compiled with XeLaTeX %%%%%
\makeatletter 
\def\beamer@framenotesbegin{% at beginning of slide
     \usebeamercolor[fg]{normal text}
     \gdef\beamer@noteitems{}% 
     \gdef\beamer@notes{}% 
}
\makeatother

%%%%% packages %%%%%
\usepackage{lipsum}					% für Testtexte
\usepackage{graphicx}				% für Bilder
\usepackage{geometry}				% für Seitenränder
\usepackage[utf8]{inputenc}			% für Zeichencodierung
\usepackage[T1]{fontenc}			% für Silbentrennung
\usepackage[english,ngerman]{babel}	% für Dokumentensprache
\usepackage{amsmath}				% für colle Mathesachen
\usepackage{amssymb}				% für colle Mathesachen
\usepackage{listings}				% für Programmiercode
\usepackage{tabularx}				% für bessere Tabellen
%\usepackage{subcaption}			% für Untertitel von Unterbildern
\usepackage{ulem}					% für Unterstreichungen usw.
%\usepackage{pifont}				% für tolle Symbole (Zapf Dingbats)
%\usepackage{enumerate}				% für bessere Aufzählungen
%\usepackage[section]{placeins}		% für Floats (setzt automatisch \FloatBarrier vor jeder neuen Section)
%\usepackage{tocloft}				% für Inhaltverzeichniseintrage ohne Nummerierung und Existenz im Text
\usepackage{pgfpages}				% für Dual-Screen-Präsentationen
\usepackage{hyperref}

%%%%% presentation template %%%%%
\setbeamertemplate{footline}[text line]{\parbox{\linewidth}{\vspace*{-8pt}Raumzeitliche Analyse von Twitter Daten\hfill\insertshortauthor\hfill\insertpagenumber}}
\setbeamertemplate{navigation symbols}{}

%%%%% presentation options %%%%%
%\setbeameroption{show notes on second screen = right}	% render note page to the right of frame
\hypersetup{pdfstartview={Fit}} 						% fits the presentation to the window when first displayed
\graphicspath{{pictures/}}								% picture path

%%%%% presentation theme %%%%%
\usetheme{Berkeley}
\usecolortheme{default}

%%%%% presentation logo %%%%%
%\logo{\includegraphics[height=1.9cm]{picture}}

%%%%% presentation informations %%%%%
\title{Raumzeitliche Analyse von Twitter Daten}
\subtitle{Softwareentwicklungprojekt I}
\author{Niklas Baumbach, Felix Juch und Martin Immel}
\date{}

%%%%% listing style %%%%%
\lstdefinestyle{mystyle}{
	basicstyle = \scriptsize\ttfamily,
	breaklines = true,
	postbreak = \text{$\hookrightarrow$},
	numbers = left,
	numbersep = 10pt,
	numberstyle = \color{Gray},
	commentstyle = \color{DarkSlateGrey},
	stringstyle = \color{Purple},
	otherkeywords = {},
	keywords = [2]{},
	keywordstyle = {\color{Blue}},
	keywordstyle = [2]{\color{OliveGreen}},
	xleftmargin = .05\textwidth,
	literate={ä}{{\"a}}1 {ö}{{\"o}}1 {ü}{{\"u}}1 {ß}{{\ss}}1 {Ü}{{\"U}}1 {Ä}{{\"A}}1 {Ö}{{\"O}}1,
	escapeinside={[*}{*]}
}


%%%%% main %%%%%
\begin{document}
	\begin{frame}
	 	\titlepage
	\end{frame}	
	
	
	\section{Inhalt des Projekts}
	\begin{frame}{Inhalt des Projekts}{}
		\textbf{Inhalte:}\\
		\begin{itemize}
			\item raumzeitliche Auswertung von Twitter-Daten zur Detektion und Analyse von Naturkatastrophen
			\item Clustern der geo-lokalisierten Tweets mit Krisenbezug
			\item Implementierung eines geeigneten Clustering-Algorithmus $\rightarrow$ FDCA (\textit{Fast Density Clustering Algorithm})
			\item Implementierung eines Frameworks zum Test des Algorithmus und zur Simulation
		\end{itemize}
	\end{frame}
	
	\begin{frame}{Inhalt des Projekts}{Clustering}
		\textbf{Clustering:}
		\begin{itemize}
			\item Bestimmung der Elementgruppen eines Datensatzes - Cluster
			\item Einteilung der Elemente in Gruppen mit ähnlichen Eigenschaften - Clustering
			\item Zugehörigkeit wird über die Bestimmung eines mehrdimensionalen Distanzma\ss es festgelegt
		\end{itemize}\bigskip
	\end{frame}
	
	\begin{frame}{Inhalt des Projekts}{FDCA}
		\textbf{FDCA:}
		
		\begin{itemize}
			\item funktioniert für numerische und kategorische Daten
			\item errechnet die Dichte der Datenpunkte
			\item bestimmen der kleinsten Distanz zu einem Datenpunkt höherer Dichte
			\item legt Clusterzentren fest (hohe Dichte, gro"se Distanz zu Datenpunkt höherer Dichte)
			\item sortiert Noise-Punkte aus (geringe Dichte, gro"se Distanz zu Datenpunkt höherer Dichte)
		\end{itemize}
	\end{frame}
	
	
	\section{Stakeholder}
	\begin{frame}{Stakeholder}{}
		\begin{columns}
			\begin{column}{.5\textwidth}
				\center \includegraphics[width=.9\textwidth]{dlr}
			\end{column}
			\begin{column}{.5\textwidth}
				\center \includegraphics[width=.8\textwidth]{fsu}
			\end{column}
		\end{columns}\bigskip
		\textbf{Betreuer:} Jens Kersten und Friedericke Klan
	\end{frame}
	
	
	% Überblick über das finale Ergebnis
	\section{}
	
	
	% Herausstellen der wesentlichen Leistungen
	\section{}
	

	% Überblick über nicht erreichte Ziele
	\section{}
	
	
	% ggf. Präsentation des Systems
	\section{}
	
	
	% Investierte Zeit pro Woche
	\section{}	
	
	
	\section{}
	\begin{frame}{Noch Fragen?}
		\includegraphics[width=\textwidth]{fragen}
	\end{frame}
\end{document}

%%%%% empty frame with notes %%%%%

%	\begin{frame}{}{}
%		\note{
%			\begin{itemize}
%				\item
%			\end{itemize}
%		}
%		
%		
%	\end{frame}